\subsection{Three constants and one transient state}

\begin{proposition}\label{C:constant-core-one-transient-H2}
Let $S=U\sqcup\{t\}$ be a finite $(N,M)$-biset.  Suppose the
right $M$-image consists of the identity and all constant maps
with values in $U$.  Suppose the left $N$-image consists of the
identity and the maps $r_j$, for $j$ in a nonempty subset
$J\subseteq U$, where $r_j$ fixes $U$ pointwise and sends $t$ to
$j$.  Let $T$ be the actual $N$-orbit set: it identifies $t$ and
all of $J$, and leaves the other points of $U$ distinct.
The actual functor from the biset collage to the ordinary
$M$-act cone on $T$ induces an isomorphism on rational second
homology, and hence on rational second cohomology.
\end{proposition}
\begin{proof}
Put $A=\ker(\mathbb QU\to\mathbb Q)$ and
$L=\ker(\mathbb QS\to\mathbb Q)$.
The identity-versus-constant character is
$\chi:M\to\{0,1\}$, and the identity-versus-retraction
character is $\psi:N\to\{0,1\}$.
The whole right action on $L$ is scalar $\chi$.
As left $N$-modules there is an exact sequence
\begin{equation}\label{C:constant-core-vector-extension}
 0\longrightarrow A\longrightarrow L
 \longrightarrow\mathbb Q_\psi\longrightarrow0,
\end{equation}
where $A$ has trivial $N$-action.

Write $X_U\subseteq X_S$ for the two collage spaces.
The coefficient extension for this pair is the pushout of
\eqref{C:constant-core-vector-extension}, tensored with
$\mathbb Q_\chi$, along the inclusion
$A\otimes\mathbb Q_\chi\to\mathbb QU$.
Consequently its relative chains are the suspended tensor of
the two action bars for $\mathbb Q_\psi$ and $\mathbb Q_\chi$.
Both bars have zero degree-zero homology, because each character
takes the value zero.  It follows that
\[
 H_2(X_S,X_U;\mathbb Q)=0,\qquad
 H_3(X_S,X_U;\mathbb Q)=
 \operatorname{Tor}^{\mathbb QN}_1(\mathbb Q,\mathbb Q_\psi)
 \otimes
 \operatorname{Tor}^{\mathbb QM}_1(\mathbb Q_\chi,\mathbb Q).
\]
The left action on $U$ is trivial.  By
Lemma~\ref{C:trivial-biset-side-H2}, collapsing the $N$-corner
identifies its second homology with that of the ordinary core
act cone.  The augmentation sequence of $\mathbb QU$ therefore
identifies
\[
 H_2(X_U;\mathbb Q)=
 A\otimes
 \operatorname{Tor}^{\mathbb QM}_1(\mathbb Q_\chi,\mathbb Q).
\]
Under these identifications the boundary of the relative pair
is $\pm\delta\otimes\id$, where
\[
 \delta:
 \operatorname{Tor}^{\mathbb QN}_1(\mathbb Q,\mathbb Q_\psi)
 \longrightarrow A
\]
is the connecting map of
\eqref{C:constant-core-vector-extension}.
This is the same literal off-diagonal-face calculation as in
Lemma~\ref{C:three-state-residual-H2zero}, now with values in
$A$: choose $u_0\in J$ and split
$\mathbb QS=\mathbb QU\oplus\mathbb Q([t]-[u_0])$;
the only off-diagonal action faces are the left extension faces;
the right faces are scalar $\chi$, and both vertex augmentations
annihilate the chosen augmentation lifts.

Taking $N$-coinvariants in
\eqref{C:constant-core-vector-extension} shows
\[
 \operatorname{im}\delta
  =\ker(A\longrightarrow L_N)
  =\ker(\mathbb QJ\longrightarrow\mathbb Q).
\]
Indeed the coinvariant relations identify $[t]$ with each $[j]$,
$j\in J$, and impose no other relations between the basis points.
Thus $L_N$ is exactly $\ker(\mathbb QT\to\mathbb Q)$.
The pair sequence now gives
\[
 H_2(X_S;\mathbb Q)=
 \bigl(A/\ker(\mathbb QJ\to\mathbb Q)\bigr)
 \otimes
 \operatorname{Tor}^{\mathbb QM}_1(\mathbb Q_\chi,\mathbb Q).
\]
The actual orbit-collapse functor induces the displayed quotient
map on the first factor.  The same augmentation calculation for
the ordinary $M$-act $T$ gives the identical expression, so this
functor is the asserted second-homology isomorphism.
\end{proof}

\begin{corollary}\label{C:three-constant-four-state-cups}
Under the hypotheses of
Proposition~\ref{C:constant-core-one-transient-H2}, if
$|U|=3$ and $|J|\geq2$, every rational degree-two cup product
on the biset collage is zero.  In particular it cannot have
type $\mathbb CP^2$.
\end{corollary}
\begin{proof}
Here $|T|=|U|-|J|+1\leq2$.
Every degree-two class pulls back from the ordinary cone on
$T$.  Theorem~\ref{C:two-cross-arrows-universal} gives zero
products when $|T|=2$; when $|T|=1$ that cone is contractible.
Naturality of cup products proves the claim.
\end{proof}
